% !TeX program = lualatex
% Compile twice: lualatex open_problems_500.tex
% All references and prose are embedded; no BibTeX database or external inputs.
\documentclass[11pt,a4paper]{article}
\usepackage[margin=24mm,headheight=14pt]{geometry}
\usepackage{fontspec}
\setmainfont{Latin Modern Roman}
\setsansfont{Latin Modern Sans}
\setmonofont{Latin Modern Mono}
\usepackage{amsmath,amssymb,mathtools}
\usepackage{unicode-math}
\setmathfont{Latin Modern Math}
\usepackage{microtype}
\usepackage{enumitem,xurl,needspace}
\usepackage[dvipsnames]{xcolor}
\usepackage{hyperref}
\hypersetup{unicode=true,colorlinks=true,linkcolor=MidnightBlue,urlcolor=MidnightBlue,
 pdftitle={500 Mathematical Problem Statements},pdfauthor={Source-annotated compilation},bookmarksnumbered=true}
\usepackage{bookmark}
\usepackage{tocloft}
\setlength{\cftsecnumwidth}{3em}
\cftsetpnumwidth{2.5em}
\cftsetrmarg{4em}
\usepackage{titlesec}
\titleformat{\section}{\Large\bfseries\raggedright}{\thesection}{0.7em}{}
\usepackage{fancyhdr}
\pagestyle{fancy}\fancyhf{}
\fancyhead[L]{\small\sffamily Mathematical problem statements}
\fancyhead[R]{\small Source edition 22}
\fancyfoot[C]{\thepage}
\setlength{\parindent}{0pt}
\setlength{\parskip}{5pt plus 1pt}
\setlength{\emergencystretch}{3em}
\setcounter{tocdepth}{1}
\setcounter{secnumdepth}{1}
\setlist[itemize]{leftmargin=3em,itemsep=5pt,topsep=4pt,parsep=0pt}
\allowdisplaybreaks
\newcommand{\N}{\mathbb{N}}
\newcommand{\Z}{\mathbb{Z}}
\newcommand{\Q}{\mathbb{Q}}
\newcommand{\R}{\mathbb{R}}
\newcommand{\C}{\mathbb{C}}
\newcommand{\F}{\mathbb{F}}
\newcommand{\A}{\mathbb{A}}
\newcommand{\PP}{\mathbb P}
\newcommand{\E}{\mathbb{E}}
\newcommand{\eps}{\varepsilon}
\newcommand{\dd}{\,\mathrm d}
\newcommand{\sref}[2]{\hyperlink{source:#1:#2}{[S#2]}}
\newcommand{\eref}[2]{\hyperlink{extra:#1:#2}{[E#2]}}
% Turn the next switch off for a shorter reading copy. The quotations preserve
% every exactTarget field, including qualifications omitted from a short summary.
\newif\ifshowcatalogue
\showcataloguetrue
\newcommand{\cataloguescope}[1]{\ifshowcatalogue
 \paragraph{Exact scope in the supplied catalogue.}
 \begin{quote}\small #1\end{quote}\fi}
% Standalone notebook for original catalogue section 160.
% Compile twice with: lualatex 160.tex
\numberwithin{equation}{section}
\hypersetup{pdftitle={Research attempt -- problem 160},pdfauthor={Research notebook; original compilation retained}}
\fancyhead[L]{\small\sffamily Research notebook}
\fancyhead[R]{\small\sffamily Problem 160 | 27 September 2026}
\begin{document}
\setcounter{section}{159}
% ================================================================
% problemId: problem.cramer-conjecture-prime-gaps
% source releaseRank: 160; source releaseStatus: open
% exactTarget SHA-256: f90425a40398d963d0073a614063dacd5beb21971485d0392d4ff7a18b3bd441
\clearpage
\section{\texorpdfstring{Cramer Conjecture on the Limsup of Normalized Prime Gaps}{Cramer Conjecture on the Limsup of Normalized Prime Gaps}}\label{160}
\subsection{Definitions and mathematical statement}
Let $p_n$ be the $n$th prime and $g_n=p_{n+1}-p_n$. With natural logarithms, the selected assertion is the exact constant statement
\[
 \limsup_{n\to\infty}\frac{g_n}{(\log p_n)^2}=1.
\]
Equivalently, for every $\eps>0$, all sufficiently large $n$ satisfy
$g_n\le(1+\eps)(\log p_n)^2$, while infinitely many satisfy
$g_n\ge(1-\eps)(\log p_n)^2$. The equality of the limsup is materially stronger than an unspecified upper bound $g_n=O((\log p_n)^2)$. This section records the catalogue's precise constant-one version; the title alone should not be used to interchange different prime-gap conjectures or heuristics.
\par\smallskip\textit{Formulation and concept references:} \sref{160}{1}.
\cataloguescope{Prove that lim sup\_\{n -> infinity\} (p\_\{n+1\} - p\_n) / (log p\_n)\textasciicircum{}2 = 1, where p\_n is the n-th prime.}
\subsection{Short English statement}
Is the largest recurring scale of consecutive prime gaps exactly the square of the logarithm of the smaller prime, with limiting upper constant one?
% BEGIN RESEARCH ADDITION ATTEMPT-160
\subsection{Research attempt}

\paragraph{Current frontier.}
The constant $1$ is the decisive issue. Granville's small-prime-conditioned
model instead suggests recurring gaps with normalized size at least
$2e^{-\gamma}>1$; this is a heuristic, not a theorem about the actual
primes \eref{160}{1}, pp.~23--24. Established landmarks include the
Baker--Harman--Pintz prime-in-short-interval estimate and the
Ford--Green--Konyagin--Maynard--Tao large-gap lower bound
\eref{160}{2}, \eref{160}{3}. Neither reaches the required logarithmic-square
scale with a positive limiting constant. A currently accessible manuscript
entitled \emph{Improved long gaps between primes} claims a stronger
iterated-logarithm lower bound \eref{160}{4}; it is not independently
verified here, and even its stated bound is $o((\log X)^2)$. It must not be
reported as a resolution of the exact constant-one target.

\paragraph{Attack route.}
Three routes are natural: construct covering congruences giving longer
prime-free intervals; estimate correlations of many prime translates;
or transfer the extreme-value calculation from a random model. The first
must produce interval length comparable to the square of the logarithm
of the constructed location. The third requires a genuine arithmetic
transfer theorem, not just numerical agreement. I pursue the second,
retaining the alternating cancellations that are normally lost in a
termwise error bound. The question becomes how accurately high factorial
moments determine an extremely rare empty window.

\paragraph{Attempt: exact finite identities.}
Let $X\ge2$ and $h\ge1$ be integers, and put
\[
 \mathcal I_X=\{X,X+1,\ldots,3X\},\qquad N_X=2X+1,
 \qquad A_m=\pi(m+h)-\pi(m).
\]
Define
\begin{align}
 E(X,h)&=\#\{m\in\mathcal I_X:A_m=0\},\label{160:empty}\\
 F_j(X,h)&=\sum_{m\in\mathcal I_X}\binom{A_m}{j}
 =\sum_{1\le a_1<\cdots<a_j\le h}
   \#\{m\in\mathcal I_X:m+a_1,\ldots,m+a_j\text{ prime}\}.
 \label{160:factorial}
\end{align}
Here $F_0=N_X$. The second equality simply counts pairs consisting of a
starting point and $j$ prime locations. It is exact, including the
inadmissible patterns whose contribution is zero.

Set $B_M=\sum_{j=0}^M(-1)^jF_j$. For each nonnegative integer $a$,
\[
 \sum_{j=0}^M(-1)^j\binom aj=
 \begin{cases}
 1,&a=0,\\
 (-1)^M\binom{a-1}{M},&a\ge1,
 \end{cases}
\]
with binomial coefficients zero beyond their usual range. This follows
from Pascal's identity by telescoping. Therefore
\begin{equation}\label{160:bonferroni}
 B_{2r+1}\le E(X,h)\le B_{2r}.
\end{equation}
Thus a short list of \emph{fixed-order} correlations is not the relevant
object: the truncation order will have to grow with $X$.

There is also an exact relation to consecutive prime gaps:
\begin{equation}\label{160:gap-identity}
 E(X,h)=\sum_{p<q\ {\rm consecutive}}
 \max\{0,\min(3X,q-h-1)-\max(X,p)+1\}.
\end{equation}
Indeed, $(m,m+h]$ contains no prime precisely when the preceding prime
$p\le m$ and the next prime $q$ satisfy $m+h<q$. Each admissible $m$ is
counted exactly once.

\paragraph{Attempt: a quantitative rare-window criterion.}
For fixed $c>0$ take $h=\lfloor c(\log X)^2\rfloor$ and use the varying
local intensity
\[
 \lambda_m=\int_m^{m+h}\frac{\dd t}{\log t},\qquad
 Q(X,h)=\sum_{m\in\mathcal I_X}e^{-\lambda_m}.
\]
Uniformly for these $m$,
$\lambda_m=c\log X+O_c(1)$, whence
\begin{equation}\label{160:poisson-scale}
 Q(X,h)\asymp_c X^{1-c}.
\end{equation}
The integral is a comparison model, not an asserted prime-counting formula.
It avoids the spurious first-moment discrepancy introduced by freezing
$\log m$ at $\log X$ across an entire interval of length $2X$.

Write
\begin{align}
 P_M&=\sum_{m\in\mathcal I_X}\sum_{j=0}^M
             \frac{(-\lambda_m)^j}{j!},\\
 \mathcal D_M&=B_M-P_M
 =\sum_{j=0}^M(-1)^j
   \left(F_j-\frac1{j!}\sum_m\lambda_m^j\right).
 \label{160:signed-discrepancy}
\end{align}
Taylor's theorem gives the rigorous error bound
\begin{equation}\label{160:taylor}
 |P_M-Q(X,h)|\le
 N_X\frac{\lambda_{\max}^{M+1}}{(M+1)!}.
\end{equation}
Choose a constant $A=A(c)$ so large that
$A\log(A/(ec))>3$, and take $M=A\log X+O(1)$ with either prescribed
parity. Using $n!\ge(n/e)^n$ and
$\lambda_{\max}=c\log X+O_c(1)$, the right side of
\eqref{160:taylor} is $o(1)$.

The following two arithmetic estimates would imply the printed target:
\begin{equation}\label{160:missing}
 \begin{aligned}
 c>1:&\quad \mathcal D_M=o((\log X)^2)
       &&\text{for the even choice of }M,\\
 0<c<1:&\quad \mathcal D_M=o(Q(X,h))
       &&\text{for the odd choice of }M.
 \end{aligned}
\end{equation}
These are \textbf{unproved hypotheses}, not consequences of the preceding
identities. Notice that they control the signed discrepancy itself, not
the sum of absolute errors of its terms.

Here is the full implication, including the uniform upper bound. For
$c>1$, \eqref{160:bonferroni}--\eqref{160:missing} give
$E(X,h)=o((\log X)^2)$. If a prime $p\in[X,2X]$ had next-prime gap
$g>h+\delta(\log X)^2+2$, then the first
$\lfloor\delta(\log X)^2\rfloor$ starts after $p$ would be empty windows.
For large $X$ all lie in $[X,3X]$, contradicting that estimate. Given
$\varepsilon>0$, take $c=1+\varepsilon/3$ and
$\delta=\varepsilon/3$. Applying the conclusion on successive dyadic
intervals proves $g\le(1+\varepsilon)(\log p)^2$ for every sufficiently
large prime $p$.

For $0<c<1$, the odd bound gives
$E(X,h)\ge(1+o(1))Q(X,h)>0$. Choose an empty window and its preceding
prime $p\le3X$. The next-prime gap is greater than $h$, and hence its
normalized size is at least
$h/(\log(3X))^2=c+o(1)$. These preceding primes tend to infinity as
$X\to\infty$. Letting $c$ approach $1$ from below proves the required
lower limsup. This proves the conditional implication, not
\eqref{160:missing}.

\paragraph{Stress test: where the proposed transfer breaks.}
Replacing \eqref{160:signed-discrepancy} by the triangle inequality is
catastrophic. Terms near $j\approx\lambda_m$ have size roughly
$e^{\lambda_m}$ before multiplication by $N_X$, whereas their alternating
sum has size $e^{-\lambda_m}$. Relative error estimates for individual
terms need not survive this cancellation. Nor does a Poisson limit for
fixed $h/\log X$ apply when that ratio is $c\log X\to\infty$.

Small-prime obstructions are material: two adjacent integers above $2$
cannot both be prime, whereas the independent model assigns a positive
pair probability. Averaging admissible singular series can correct some
fixed-order statistics, but no growing-order cancellation theorem with
\eqref{160:missing}'s accuracy has been supplied. Granville's calculation
is specifically a warning that the constant-one transfer may be false
\eref{160}{1}. Proving the first line of \eqref{160:missing} in ignorance
of that warning would effectively assume the desired upper extreme-value
law.

The direct counterexample route also fails at its scale. The elementary
composites $N!+2,\ldots,N!+N$ certify only a length of order $N$ at a
location of logarithm of order $N\log N$. That certified lower bound,
divided by the squared logarithm, tends to zero; this statement is about
the certificate, not an upper bound for the actual surrounding prime gap.
The accompanying script checks \eqref{160:bonferroni} and
\eqref{160:gap-identity} at $X=100,h=24$. These are finite consistency
checks, not evidence establishing either limit in \eqref{160:missing}.

\paragraph{Outcome.}
\textbf{Outcome: Exploratory approach with unresolved gap.}

The exact empty-window identities and the implication from
\eqref{160:missing} are proved. The attempt isolates the needed
cancellation scale and shows why fixed-order moments, mean density, and
naive factorial constructions do not reach it. Neither the constant-one
upper bound nor infinitely many actual prime gaps exceeding that constant
has been established. The competing arithmetic heuristic remains a reason
to investigate disproof rather than to presume the catalogue's equality.

% Keep the local source heading and initial references together.
\Needspace{12\baselineskip}
% END RESEARCH ADDITION ATTEMPT-160
\subsection{Sources}
\begin{itemize}\raggedright
\item[\textnormal{[S1]}]\hypertarget{source:160:1}{} Weisstein, E. W., Cramer Conjecture, From MathWorld - A Wolfram Resource, updated 16 August 2026.
\url{https://mathworld.wolfram.com/CramerConjecture.html}.
{\small\textit{Catalogue formulation source.}}
% BEGIN RESEARCH ADDITION SOURCES-160
\item[\textnormal{[E1]}]\hypertarget{extra:160:1}{} Andrew Granville, ``Harald Cram\'er and the distribution of prime numbers,'' \emph{Scandinavian Actuarial Journal} 1995, no. 1, 12--28, especially pp. 23--24.
\url{https://dms.umontreal.ca/~andrew/PDF/cramer.pdf}.
{\small\textit{Small-prime-conditioned heuristic; not a proof of a prime-gap limsup.}}
\item[\textnormal{[E2]}]\hypertarget{extra:160:2}{} R. C. Baker, G. Harman and J. Pintz, ``The difference between consecutive primes, II,'' \emph{Proc. London Math. Soc.} (3) 83 (2001), 532--562. Main prime-in-short-interval result.
\url{https://doi.org/10.1112/plms/83.3.532}.
{\small\textit{Established polynomial-scale upper-bound landmark, not a logarithmic-square estimate.}}
\item[\textnormal{[E3]}]\hypertarget{extra:160:3}{} Kevin Ford, Ben Green, Sergei Konyagin, James Maynard and Terence Tao, ``Long gaps between primes,'' \emph{J. Amer. Math. Soc.} 31 (2018), 65--105; arXiv:1412.5029. Theorem 1.
\url{https://arxiv.org/abs/1412.5029}.
{\small\textit{Established large-gap lower bound.}}
\item[\textnormal{[E4]}]\hypertarget{extra:160:4}{} OpenAI, ``Improved long gaps between primes,'' eight-page online manuscript, introductory theorem; accessed 27 September 2026. The consulted PDF does not supply a journal publication record.
\url{https://cdn.openai.com/pdf/51126fac-1b68-4128-9666-c908bcc16033/long_gaps.pdf}.
{\small\textit{Current proof claim used only to compare the stated scale with the target; not independently verified or endorsed.}}
% END RESEARCH ADDITION SOURCES-160
\end{itemize}

\end{document}
